\documentclass[11pt,a4paper]{article}

% Self-contained Overleaf source; no external bibliography or figures required.
\usepackage[margin=1in]{geometry}
\usepackage{graphicx}
\usepackage{booktabs}
\usepackage{amsmath}
\usepackage[T1]{fontenc}
\usepackage{authblk}
\usepackage{tabularx}
\usepackage{xurl}
\usepackage{hyperref}
\hypersetup{hidelinks,
  pdftitle={Auditing the OPAS Orbital Proximity Alert System's Safe-Window Scanner},
  pdfauthor={Nadim Baboun}}

\title{Auditing the OPAS Orbital Proximity Alert System's Safe-Window Scanner}
\author[1]{Nadim Baboun}
\affil[1]{Al-Quds University}
\date{7 October 2026}

\begin{document}

\maketitle

\begin{abstract}
The Orbital Proximity Alert System (OPAS) estimates candidate launch windows
from two-line element (TLE) sets and SGP4 propagation. This study audits its
\texttt{/safe-windows} scanner against an independently implemented dense
numerical reference. In a frozen six-hour launch horizon, every reference
launch sample covered by the original scanner's reported ISS, Starlink and
sun-synchronous windows was obstructed at the scanner's own screening
radius; a geostationary control case had no obstructed reference samples.
An initial four-fix scanner reduced this sampled false-safe rate to 9.6\%
for its ISS window and 21.1\% for its Starlink window, with all ten
investigated residual object--launch pairs traced to finite waypoint
spacing. A follow-up scanner refines close approaches between waypoints,
uses exact clock launch grids, and verifies candidate windows at 10-second
spacing. Its six-hour ISS search agreed with the independent reference on
all 691 checked launch times, including zero observed false-safe samples
among 565 launches classified clear. Finer verification exposed 52 unsafe
launch samples inside apparently clear minute-grid windows, leaving no
qualifying window of at least 15 minutes. Archived ISS, Starlink and SSO
pair replay also produced zero disagreements. These results support the
sampling-gap diagnosis and demonstrate improved sampled clearance within
OPAS's frozen post-ascent model. They do not certify continuous launch-time
clearance or real-world launch safety. Anisotropic screening, calibrated
orbital uncertainty, ascent coverage and further runtime work remain open.
\end{abstract}

\section{Introduction}

The Orbital Proximity Alert System (OPAS) is a lightweight web tool for
screening planned satellite and rocket trajectories for proximity to
tracked objects\footnote{Source repository:
\url{https://github.com/nadeemtsf/OPAS}. The follow-up measurements belong to the
\texttt{fix/refine-close-approaches} branch, not the repository's default
branch.}. Given a launch site,
target altitude and inclination, its \texttt{/safe-windows} endpoint
returns candidate launch-time intervals that its model classifies as
clear. OPAS ingests TLE sets from Space-Track, propagates object positions
using SGP4 via Skyfield, and screens a simplified vehicle trajectory
against objects selected near the target altitude.

OPAS was developed through rapid, AI-assisted prototyping. Its claimed
safe windows need independent verification, particularly where sampling
reduces web-request cost. This report examines missed close approaches
under its own model and threshold, their causes and subsequent corrections.

Throughout this report, ``clear'' means no reference-resolved encounter
inside the chosen radius for a particular checked launch time and modeled
flight scope. Consecutive clear launch samples are not, by themselves, a
certificate that every intervening launch time is clear. Historical
measurements and follow-up results are kept separate to preserve their
different code versions, launch grids and validation coverage.

\section{Background}
\label{sec:background}

\subsection{TLE/SGP4 propagation accuracy}

TLE propagation accuracy is object- and epoch-dependent. Priestley and
Handley report an often-quoted heuristic of approximately 1~km error at
epoch, growing by 1--3~km per day; this is context for numerical precision,
not a calibrated uncertainty law for every object~\cite{jaxsgp4}. Rich's
LEO study, which included the Hubble Space Telescope, reported
SGP4-propagated position error within 1~km over seven days in the cases
examined~\cite{afit_sgp4}. These observations should not be converted into
a universal TLE-age penalty. A screening margin needs evidence appropriate
to the objects, tracking quality and mission being protected.

\subsection{Documented conjunction-screening practice}

NASA's 2023 handbook assigns screening volumes according to orbital and
mission characteristics and distinguishes conjunction screening from risk
assessment~\cite{nasa_handbook}. Hejduk and Pachura's six-month study used
an intentionally large 50~km radial by 250~km in-track by 250~km cross-track
volume to investigate screening performance. Their approximately 700~km
orbit group used a smaller ellipsoid with \emph{semi-axes} of 0.5, 17 and
20~km in those directions~\cite{nasa_screening}; the handbook also reports
a 0.5~km~$\times$~17~km~$\times$~20~km screening example~\cite{nasa_handbook}.
The 0.5~km value is a radial dimension, not a spherical collision radius.

These documented examples are mission-specific screening volumes rather
than universal final safety thresholds. The handbook also advises against
using public TLEs alone for conjunction risk assessment because their
accuracy is insufficient for that purpose~\cite{nasa_handbook}. OPAS's
numerical audit consequently concerns a model-based proximity screen,
not an operational collision-avoidance certification. The historical
examples cited here are identified by publication version; no claim is
made that they specify every mission's present screening configuration.

\section{Methodology}
\label{sec:methodology}

\subsection{Frozen data and scanner versions}

The MongoDB debris collection was exported once to a checksummed snapshot
of 32,517 objects. Its epoch, $T_0$, is
\texttt{2026-09-21T08:52:54.811735+00:00}. The requested six-hour
\emph{launch} horizon ends at
\texttt{2026-09-21T14:52:54.811735+00:00}; each launch is then screened
over its own modeled post-ascent flight. TLE ages in retained measurements
are evaluated relative to $T_0$, independent of when a script is run.
The snapshot identifier and result paths are recorded in
Appendix~\ref{sec:records}.

Three scanner stages are distinguished:
\begin{enumerate}
  \item \textbf{Original scanner:} unmodified OPAS commit
  \texttt{a89f2e72}, used for the baseline audit and Table~\ref{tab:baseline}.
  \item \textbf{Initial four-fix scanner:}
  the research scanner, with corrected frozen-age results
  dated 28 September 2026. Its production counterpart is
  \texttt{fix/scanner-accuracy} at \texttt{50fb28f}. Its historical results
  appear in Table~\ref{tab:fixed}.
  \item \textbf{Follow-up scanner:}
  \texttt{fix/refine-close-approaches}, based on \texttt{50fb28f}.
  The full numerical search measured local commit \texttt{19548cd}.
  Subsequent integration fixes are in local commit \texttt{9637985};
  their separate verification is described below~\cite{opas_validation}.
\end{enumerate}

The original Windows study used Python 3.13.14, NumPy 2.4.4, Skyfield 1.54
and the optional C++ extension. The Linux follow-up used Python 3.12.14,
NumPy 2.3.5, Skyfield 1.55 and \texttt{sgp4} 2.27; timings are not a
controlled before-and-after comparison.

Research scanner and sweep outputs with unfrozen TLE age were invalidated
and archived. Only corrected records are used; the original
baseline/reference pipeline was unaffected.

\subsection{Baseline measurement}

The original \texttt{scan\_windows} was run over the launch horizon for
ISS (420~km, 51.6\textdegree{} inclination), Starlink (550~km,
53.0\textdegree{}), SSO (705~km, 98.2\textdegree{}), polar orbit
(800~km, 90\textdegree{}) and the preset named ``GEO transfer''
(35,786~km, 6.0\textdegree{}). The last preset still uses OPAS's simplified
trajectory model; its name does not imply a physically modeled transfer
trajectory. Each run recorded the actual windows returned to a user.

\subsection{Dense reference implementation}

The archived reference shares OPAS's snapshot, SGP4 propagator and vehicle
generator. It resolves the modeled flight at approximately one-second
spacing and refines minima by parabolic interpolation of squared distance
around the nearest sample. This refinement is exact for locally
straight-line relative motion and an approximation for curved motion.
It supplies a validated numerical reference rather than physical ground
truth.

Two planted-object tests recovered known minima to within 0.02~km after
refinement. A one-second versus two-second comparison over 38,000
object--launch pairs found a median distance disagreement of 0.5~m and
zero classification flips at 5, 10, 25 or 50~km radii. These checks support
the reference resolution in the tested cases; they are not an exhaustive
error bound for every pair or launch time.

Archived launch offsets derive from numerical flight spacing and differ
slightly from clock minutes. Historical cases retain their recorded offsets;
the follow-up compares identical exact-clock timestamps.

\subsection{Validating the reported windows}

For each original window, the reference classifies covered approximately
minute-spaced launch samples using OPAS's own per-object radius. The
\emph{sampled false-safe rate} is the number of covered launch samples
inside that reported window which the reference marks obstructed, divided
by the number of covered reference launch samples. It is not the fraction
of continuous elapsed time proved unsafe. Missing endpoint coverage is
reported separately and is not treated as clearance. Windows are also
clipped to the requested horizon to separate the boundary defect from
sampling errors.

A sweep varied coarse step, fine step and inner-orbit sample count. Its
parameterized scanner first reproduced the original ISS, Starlink and SSO
windows and evaluation counts, testing sampling under the original model.

\subsection{Initial four-fix scanner}
\label{sec:initial-method}

The first correction (i) increased trajectory generation to 600 intervals
per modeled orbit, producing 601 waypoints before post-ascent selection,
and evaluated every retained waypoint; (ii) replaced the 50~km base and
age inflation capped at 200~km with a 10~km LEO base and capped linear
age penalty; (iii) tightened launch verification from two-minute to
one-minute spacing; and (iv) clamped returned windows to the requested
horizon. The LEO radius used in the retained data is
\begin{equation}
  r_i = 10 + \min\!\left(0.5\max(a_i,0),\,5\right)\quad\mathrm{km},
  \label{eq:radius}
\end{equation}
where $a_i$ is TLE age in days. Higher-orbit presets retain a 50~km base.
This is an engineering screening choice, not a measured uncertainty model.

Historical residual pairs were replayed at exact recorded offsets and
radii. Consecutive clear archived launches were also enumerated independently
of window formation. Their endpoint span is $t_{\rm last}-t_{\rm first}$:
$N$ minute samples do not span $N$ minutes.

\subsection{Follow-up refinement and independent verification}
\label{sec:followup-method}

The follow-up preserves the 600-interval trajectory, radius formula,
SGP4 propagation and post-ascent scope. For each interval between retained
flight waypoints, it first screens the minimum distance along the chord
of relative Earth-fixed motion. Its screening allowance is
\begin{equation}
  \epsilon = \frac{a_{\max}(\Delta t)^2}{8} + 0.02~\mathrm{km},
  \qquad a_{\max}=0.05~\mathrm{km\,s^{-2}}.
  \label{eq:allowance}
\end{equation}
The acceleration ceiling and allowance are engineering assumptions, not
universal bounds. Potential encounters receive cubic Hermite vehicle
interpolation and fresh satellite propagation, with up to 20 golden-section
iterations. Failed or non-finite predictions cannot establish clearance.

Exact clock grids are independent of flight spacing. Ten-minute discovery
leads to minute checks and ten-second candidate verification. Unsafe samples
split runs; checked-clear endpoints must span at least 15 minutes. All
eligible candidates are considered before returning the five longest,
without extrapolating boundaries. Intervening launches remain unchecked.

The ISS test used 11,997 objects within the existing 420~km $\pm$100~km
altitude filter. Its independent reference used a regenerated 1.000079~s
vehicle grid, interpolated one-second satellite ephemerides and parabolic
distance minimization. Minima at or below radius plus 20~m were directly
repropagated. It checked 361 exact-minute and 375 finer launches, covering
every production timestamp plus 45 extra minute timestamps. Flight checking
starts at 603.356~s for production and 600.047~s for the reference because
their grids round the post-ascent boundary differently.

Supplementary replay used exact-offset ISS, Starlink and SSO pairs inside
the horizon and production scope, including clear controls within radius
plus 20~km. A direct approximately 0.2-second control checked six
threshold-nearest ISS pairs; neither is a complete finer-grid catalogue run.

Final integration fixes report unavailable orbital data explicitly and
preserve selected UTC timestamps. Tests and pair replay were rerun, but
the full search was not repeated or retimed. Interval mathematics, vehicle
generation and window discovery remain byte-identical. All 11,997 ISS
objects construct, so the new failure guards are not reached by the
snapshot~\cite{opas_validation}.

\section{Results}
\label{sec:results}

\subsection{Original scanner: false-safe samples in crowded shells}

Every covered reference launch sample inside the original ISS, Starlink
and SSO windows was obstructed; the GEO control had none
(Table~\ref{tab:baseline}).

\begin{table}[htbp]
\centering
\small
\begin{tabular}{lrrr}
\toprule
Preset & Window duration & Reference samples & Sampled false-safe rate \\
\midrule
ISS (420~km)      & 24~min & 24 & 100\% \\
Starlink (550~km) & 46~min & 46 & 100\% \\
Starlink (550~km) & 16~min & 16 & 100\% \\
SSO (705~km)      & 20~min & 20 & 100\% \\
GEO preset       & 360~min (clipped) & 360 & 0\% \\
\bottomrule
\end{tabular}
\caption{Original scanner results over the six-hour launch horizon.
Rates are fractions of covered reference launch samples, not continuous
duration estimates. Polar orbit returned no window.}
\label{tab:baseline}
\end{table}

\subsection{Root cause and boundary defect}

The original inner check used approximately a dozen intervals per orbit,
yielding roughly 14 checked points about six minutes apart. At a flat
50~km radius, every tested minute-grid launch in ISS, Starlink and SSO
had at least one reference-resolved encounter, with means of 8.8--12.2
objects per launch. Sparse waypoint checks therefore missed encounters
which occurred between their evaluation times.

Doubling the sweep's inner sample count from 12 to 24 reduced the returned
window count to zero for those three presets at every tested coarse step.
This is consistent with the sampling diagnosis under the original
radius. An independent boundary defect let refinement evaluate and return
times outside the requested horizon: the original GEO result extended
eight minutes before and after the range, yielding 376 minutes before
clipping to 360 minutes.

\subsection{Initial four-fix scanner and residual waypoint misses}

The first correction returned one ISS window and one Starlink window,
shown in Table~\ref{tab:fixed}; SSO returned none. These historical results
belong to the four-fix research scanner, not to the follow-up branch.

\begin{table}[htbp]
\centering
\small
\begin{tabular}{llrrr}
\toprule
Preset & Window (UTC) & Samples & Unsafe & Sampled false-safe rate \\
\midrule
ISS      & 10:33--11:25 (52~min) & 52 & 5 & 9.6\% \\
Starlink & 10:03--10:22 (19~min) & 19 & 4 & 21.1\% \\
\bottomrule
\end{tabular}
\caption{Initial four-fix scanner windows, graded at its exact per-object
10--15~km radius. Times are abbreviated to minutes; recorded timestamps
retain seconds and fractions. SSO returned no window.}
\label{tab:fixed}
\end{table}

The nine obstructed launches involved ten object--launch pairs. Replayed
at recorded offsets, all ten had waypoint minima outside the radius despite
reference-refined minima inside: five ISS and five Starlink pairs confirmed
between-waypoint misses. Spacing was approximately 9.28~s and 9.55~s,
allowing roughly 139--143~km of relative motion at 15~km/s. This diagnoses
the examined cases, not every possible failure mode.

\subsection{Archived clear-sample runs and the duration correction}
\label{sec:archived-runs}

The archived enumeration added one minute to each first-to-last sample
span. Its five ISS runs labelled 15--19 minutes contained 15--19 samples;
three span only about 14 minutes. Recounting archived distances and frozen
radii leaves two ISS spans of at least 15 minutes
(Table~\ref{tab:runs}); Starlink and SSO have none.

\begin{table}[htbp]
\centering
\small
\begin{tabular}{lrrr}
\toprule
Preset & Longest run: samples & Endpoint span (min) & Spans $\geq$15 min \\
\midrule
ISS      & 19 & 18.0014 & 2 \\
Starlink & 14 & 13.0003 & 0 \\
SSO      &  6 &  4.9997 & 0 \\
\bottomrule
\end{tabular}
\caption{Recount of consecutive clear samples on the archived nominal
minute launch grids, using $t_{\rm last}-t_{\rm first}$ rather than adding
one minute. The second qualifying ISS endpoint span is 16.0013 minutes.
These are sampled runs, not continuously certified clear intervals.}
\label{tab:runs}
\end{table}

These are sampled-clear candidates, not continuous certificates. The
historical first-five-candidate cap and truncation at unsafe checks are
search limitations, but their individual effects on missed runs were not
experimentally isolated. The archived grid also differs from clock minutes.

\subsection{Follow-up ISS launch classification and window verification}
\label{sec:followup-results}

The instrumented production search made 691 launch checks: 316 during
minute discovery and 375 at finer spacing. Of these, 565 were clear and
126 obstructed. The independent calculation supplied 736 launch results,
including every exact production timestamp. Table~\ref{tab:followup}
summarizes their agreement~\cite{opas_validation}.

\begin{table}[htbp]
\centering
\begin{tabularx}{\linewidth}{Xr}
\toprule
Follow-up ISS measurement & Result \\
\midrule
Production launch samples checked & 691 \\
Independent launch samples checked & 736 \\
Matched production classifications & 691/691 \\
Production launches classified clear & 565 \\
Observed false-safe launch samples & 0/565 (0\%) \\
Observed false alarms & 0 \\
Returned windows with endpoint span $\geq$15 min & 0 \\
\bottomrule
\end{tabularx}
\caption{Follow-up classification agreement at exact checked launch times
in the frozen ISS post-ascent model. The false-safe denominator is clear
launch samples, not returned windows or continuous elapsed time.}
\label{tab:followup}
\end{table}

The independent exact-minute grid produced two apparently clear endpoint
spans, 33 and 23 minutes. Ten-second checking added 280 intervening launch
samples inside those spans, of which 52 were obstructed
(18.57\%; Table~\ref{tab:minute-gaps}). These are missed unsafe samples of
\emph{minute-only launch verification}, separate from the follow-up
scanner's observed false-safe rate. The final candidate refinement split
these apparent windows and returned no qualifying 15-minute window.
Zero returned windows provides no positive returned-window case to certify,
and does not mean every launch sample was obstructed.

\begin{table}[htbp]
\centering
\small
\begin{tabular}{lrrrr}
\toprule
Minute-grid span (UTC) & Min & Extra samples & Unsafe & Unsafe fraction \\
\midrule
10:51:54--11:24:54 & 33 & 165 & 25 & 15.15\% \\
12:26:54--12:49:54 & 23 & 115 & 27 & 23.48\% \\
\midrule
Combined & 56 & 280 & 52 & 18.57\% \\
\bottomrule
\end{tabular}
\caption{Unsafe intervening launch samples inside apparent exact-minute
ISS windows. All endpoint timestamps retain the common fractional part
\texttt{.811735}; fractions are omitted here for space. These spans were
generated by the new independent minute reference, not copied from the
historical four-fix scanner output.}
\label{tab:minute-gaps}
\end{table}

These spans differ from archived runs because the grids, reference details
and boundary rounding differ; their durations are not a controlled
scanner-only comparison. Follow-up agreement compares identical timestamps.

\subsection{Supplementary replay, integration and runtime}

On the earlier integration version, archived pair replay detected all 451
selected unsafe pairs and retained all 3,089 nearby clear controls, with
zero disagreements across ISS, Starlink and SSO. Six additional nearby
pairs had archived closest approaches before the production flight scope
and were excluded. This is pair-level detector verification, not a full
Starlink or SSO window-finder rerun. The six selected approximately
0.2-second ISS controls showed zero classification changes and zero
production pair disagreements; this does not establish that the entire
catalogue was checked at that resolution.

Earlier integration checks passed: 38 backend tests in 6.789~s, nine frontend
timestamp tests across UTC, Jerusalem and New York, and frontend build
and lint. These are software checks, not physical safety validation.

The six-hour ISS search took 1,106.109~s (18.44 minutes), versus
948.989~s (15.82 minutes) for the reference. Production, using four
processes, overlapped with threaded reference work and part of an archived
replay under an eight-core CPU quota; BLAS/OpenMP used one thread.
Timings include instrumented classification phases but exclude input
setup, startup and final serialization. They describe shared load, not
isolated latency or a controlled original-scanner speed comparison.

\subsection{Later profiling and full preset comparison}

Matched six-hour searches compare the pre-optimization branch with exact
batched SGP4, shared Skyfield transforms and reused mission geometry. Radii,
altitude selection, interval refinement and strict prediction errors remain
unchanged. Profiling identified repeated orbital-position calculations as
the main cost; worker initialization and scheduling contributed seconds on
this host. A separate endpoint regression now retains near-qualifying
10-second runs for checked 5-second endpoints without reducing the required
900-second duration.

\begin{table}[htbp]
\centering
\small
\caption{Later matched searches: runtime and independent sampled agreement.}
\label{tab:performance}
\begin{tabular}{lrrrr}
\toprule
Preset & Before (s) & After (s) & Agreement & False-safe \\
\midrule
ISS (420 km) & 1292.461 & 707.157 & 691/691 & 0/565 \\
Starlink (550 km) & 708.826 & 436.498 & 418/418 & 0/281 \\
SSO (705 km) & 163.709 & 99.228 & 271/271 & 0/144 \\
\bottomrule
\end{tabular}
\end{table}

All 1,380 checked classifications are unchanged, with zero disagreements
against a separate approximately one-second flight reference. The reference
covers all 4,321 five-second launch times per preset, finding no qualifying
15-minute span in ISS, Starlink or SSO. Forty-two unfiltered full-catalogue
controls and 18 selected approximately 0.2-second pair controls also agree.
The final software suite passes 77 backend and 25 frontend tests, build and
lint. Exact-input captures and complete journals pass checksum/count audits.

The main workloads ran sequentially with four workers and one BLAS/OpenMP
thread under the same eight-core quota; small software controls overlapped
parts of the initial series, and capture reruns followed a workspace
reconnect. Times include scanner/worker preparation and capture, excluding
snapshot reading and trajectory generation. They are
single observed runs, separate from the earlier shared-load timings and the
fresh Windows run, whose exact inputs were unavailable. Agreement remains
sampled-model evidence; no positive real-catalogue window was available.

\section{Discussion}
\label{sec:discussion}

\subsection{Screening radius and uncertainty are separate from sampling}

The original 50~km sphere was OPAS's final proximity threshold, whereas
the large NASA example was a screening dataset used to study event
capture and filtering~\cite{nasa_screening}. Their similar radial scale
does not make the thresholds operationally equivalent. Moreover, the
smaller documented 0.5~km radial semi-axis belongs to an anisotropic
volume with much larger in-track and cross-track axes, not a 0.5~km
sphere~\cite{nasa_screening}.

The radius and 0.5~km/day penalty are engineering choices, not measured
catalogue errors or mission-specific collision-risk criteria. Changing
them changes obstruction. Starlink and SSO unsafe-launch fractions rose
from about 10--11\% at a flat 5~km radius to 29--45\% at 10~km. Each
scanner is graded at its own criterion, which does not validate that
threshold as actual mission risk.

An isotropic radius is also a simplification. Mission-specific
component-wise screening volumes and covariance-based risk assessment
address a different question from OPAS's spherical nominal-distance
test~\cite{nasa_handbook,nasa_screening}. The present follow-up improves
resolution of that existing test. It does not add covariance propagation,
anisotropic screening or a validated operational uncertainty model.

\subsection{Flight-time and launch-time gaps require distinct fixes}

The ten historical misses establish a between-waypoint problem during
modeled flight. Interval screening and closest-approach refinement address
that mechanism without replacing SGP4. The follow-up pair replay supports
this improvement in the investigated archived cases.

The 52 obstructed intervening launches expose a separate launch-time gap:
minute-spaced clear results do not guarantee intervening clearance. This
supports the sampling diagnosis while narrowing claims about genuinely
clear intervals. The historical waypoint diagnosis remains valid for its
cases, not an exhaustive account of window errors.

\subsection{What zero disagreements and zero windows establish}

The 691/691 agreement and 0/565 false-safe samples concern selected
timestamps in one model, not independent random trials or a universal
0\% false-safe rate.

No qualifying ISS window means apparent spans failed finer checking.
It neither validates a positive returned window nor makes every launch
unsafe: 565 checked launches were clear. Runtime remains substantial
for a synchronous request.

\section{Limitations}
\label{sec:limitations}

\textbf{Physical and endpoint scope.} This study evaluates
\texttt{/safe-windows} nominal-distance screening. The probability-of-
collision scoring and detector mathematics of the separate
\texttt{/alert} endpoint were not independently validated by these
experiments. Preserving its selected launch timestamp is an integration
fix, not validation of that endpoint's risk model.

\textbf{Vehicle model and ascent.} The vehicle follows OPAS's simplified
orbit construction, which includes a schematic initial altitude ramp.
Safe-window flight checking starts after approximately 600~s, so ascent
safety is not established. Earlier production and dense-reference starts
differ slightly because of flight-grid rounding; the later preset reference includes the exact
production start. The finer reference regenerates the nominal vehicle,
whereas production interpolates stored waypoints.
Neither represents a vehicle-specific guided launch trajectory.

\textbf{Catalogue filtering and physical uncertainty.} The existing
target-altitude filter is retained. Objects outside it, untracked debris,
TLE prediction errors, maneuvers and vehicle trajectory uncertainty are
not independently bounded by the experiment. Both numerical methods
share SGP4 and the vehicle model, so their agreement can coexist with
physical model error.

\textbf{Discrete launch times and numerical refinement.} Ten-second
launch samples leave intervening times unchecked. Flight interval
refinement relies on interpolation, finite minimization and the assumed
acceleration allowance in Equation~\ref{eq:allowance}. Its tests and
resolution controls support the investigated cases, not a universal
mathematical guarantee for arbitrary trajectories or invalid orbital data.

\textbf{Validation coverage.} All measurements use one snapshot, with
different coverage across stages. The later full-horizon references cover
ISS, Starlink and SSO; polar and GEO follow-up searches remain unreported.
None of the three later presets
yielded a qualifying positive window.

\textbf{Historical conventions.} Archived offsets drift slightly from
clock minutes. Original comparisons retain those offsets and sample-based
denominators; recounting corrects an extra-minute duration convention.
Comparison with exact-clock results requires matching timestamps.

\textbf{Search and versions.} The follow-up removes the historical
first-five-candidate restriction and splits unsafe runs, but coarse
discovery is not a continuous-time completeness proof. Earlier integration
checks remain separate from the earlier full search;
the later matched stage records its own source hashes and reference coverage.

\textbf{Threshold and runtime.} The radius and age penalty remain
engineering assumptions. Earlier timings reflect overlapping instrumented
calculations. The later matched stage measures speedup on one environment, with main workloads
sequential; single-run timings do not establish Windows or general latency.

\section{Conclusion and Future Work}
\label{sec:conclusion}

The original crowded-shell windows were obstructed at every covered
reference launch sample. Four initial fixes reduced sampled false-safe
rates to 9.6\% for ISS and 21.1\% for Starlink; all ten investigated
residual pairs confirmed waypoint-gap misses. Correcting archived duration
leaves two ISS spans of at least 15 minutes rather than five, still
without continuous-clearance certification.

The follow-up's interval refinement, exact-clock grids and ten-second
verification match all 691 checked ISS classifications, with no observed
false-safe samples among 565 clear launches. Finer checks expose 52
obstructed launches inside apparent minute-grid windows, leaving none
qualifying at 15 minutes. This supports improved sampled screening,
not continuous or physical safety.

Remaining work includes bounding intervening launch times and refinement
error, modeling vehicle-specific ascent, calibrating uncertainty and
evaluating anisotropic screening. Broader snapshots, presets and positive
returned-window cases are needed. Later matched benchmarks reduce runtime
while preserving checked classifications; broader runtime measurements
remain needed.

\appendix
\section{Evidence and Reproducibility Records}
\label{sec:records}

The report evidence index, \path{validation/EVIDENCE.md}, maps each table
and supplementary check to exact files, code versions and the snapshot
checksum. It covers the baseline comparisons, corrected four-fix results,
waypoint-gap evidence, duration recount and follow-up checks. The historical
inventory is in \path{research/results/README.md} on the research branch
at \texttt{8fef312}.

The duration recount uses archived reference distances and frozen radii
within the original horizon tolerance ($-1$ to $21603$ seconds). It reports
last minus first launch offset without rerunning propagation.

The index distinguishes the full numerical search from later integration
checks and links their manifests. Reproduction commands are in
\path{validation/README.md}.

\begin{thebibliography}{9}

\bibitem{jaxsgp4}
C. Priestley and W. Handley,
``jaxsgp4: GPU-accelerated mega-constellation propagation with batch parallelism,''
\emph{arXiv preprint arXiv:2603.27830}, 2026.
\url{https://arxiv.org/abs/2603.27830}.

\bibitem{afit_sgp4}
A. T. Rich,
``Investigating Analytical and Numerical Methods to Predict Satellite Orbits Using Two-Line Element Sets,''
M.S. thesis, Air Force Institute of Technology, AFIT-ENY-MS-17-M-286, 2017.
\url{https://scholar.afit.edu/etd/1722/}.

\bibitem{nasa_screening}
M. D. Hejduk and D. A. Pachura,
``Conjunction assessment screening volume sizing and event filtering in light of natural conjunction event development behaviors,''
\emph{AAS/AIAA Astrodynamics Specialist Conference}, Stevenson, WA,
August 2017. NASA Technical Reports Server document 20170007928.
\url{https://ntrs.nasa.gov/citations/20170007928}.

\bibitem{nasa_handbook}
NASA Office of the Chief Engineer,
``NASA Spacecraft Conjunction Assessment and Collision Avoidance Best Practices Handbook,''
NASA/SP-20230002470, Rev.~1, February 2023.
\url{https://ntrs.nasa.gov/citations/20230002470}.

\bibitem{opas_validation}
OPAS project,
``Focused scanner verification and final integration records,''
6 October 2026. Local measured commit \texttt{19548cd} and final integration
commit \texttt{9637985}; evidence index and reproduction guide identified
in Appendix~\ref{sec:records}. These records had not been pushed or published
at the time of this report.

\end{thebibliography}

\end{document}
